\documentclass[letterpaper]{article}
\usepackage[submission]{aaai2027}
\usepackage[hyphens]{url}
\usepackage{graphicx}
\urlstyle{rm}
\def\UrlFont{\rm}
\usepackage{natbib}
\usepackage{caption}
\frenchspacing

\usepackage{amsmath}
\usepackage{amssymb}
\usepackage{booktabs}
\usepackage{array}

\pdfinfo{/TemplateVersion (2027.1)}
\setcounter{secnumdepth}{0}

\title{When Learning Paths Lie:\\Gradient Retention under Hidden Relations}
\author{Anonymous Submission}
\affiliations{}

\begin{document}
\maketitle

\begin{abstract}
Little movement in policy space admits two incompatible explanations: the policy has converged, or the learning signal has disappeared.
Path-length diagnostics cannot separate them.
We study one source of this ambiguity in multi-agent reinforcement learning: \emph{relational aliasing}, where states requiring different behavior share an observation because teammate identity, partner role, or task assignment is hidden.
The resulting relation-conditioned policy gradients can cancel when aggregated.
We call this failure \textbf{Information-Induced Gradient Contraction} (IIGC).

For relation-conditioned gradients $g_A$ and $g_B$, we define total gradient energy
$\mathcal{E}=(\|g_A\|^2+\|g_B\|^2)/2$ and the \textbf{Relational Gradient Retention Ratio}
$\kappa=\|(g_A+g_B)/2\|^2/\mathcal{E}$.
An exact decomposition shows that $\kappa$ is the fraction of relation-conditioned gradient energy retained after mixing.
This interpretation also exposes a necessary boundary: $\kappa$ is undefined when $\mathcal{E}=0$, and $\kappa=1/2$ can mean that only one relation supplies gradient energy.
We therefore pair behavioral path length with an energy gate and $\kappa$, rather than treating any scalar as a stand-alone convergence test.

Across 22 PPO runs in the 510K card game, hidden teams produce the shortest behavioral paths, whereas a controlled Toy task produces a $44\times$ longer path under hidden context despite chance return.
Toy gradients resolve the contradiction: $\kappa$ falls from $0.726$ to $0.039$ over ten seeds.
Overcooked supplies a stress test in which zero return makes gradient energy vanish, demonstrating why the gate is essential.
Measurements with policy-gradient and TD estimators further show that retention is estimator-specific, not a universal algorithm score.
\end{abstract}


\section{Introduction}

A policy that barely changes can be finished or unable to start.
Both cases yield a short trajectory through parameter or behavior space, but they call for opposite decisions: stop training in the first case; repair the information structure in the second.
Return eventually separates them, but only after the run has spent the compute that a diagnostic is meant to save.
The question is therefore not whether a policy moved, but whether the update rule had a coherent direction in which to move it.

This distinction surfaced in a four-player card game.
Across 22 PPO runs in 510K, the behavioral path integral decreased from solo play to fixed teams and was shortest when team identity was hidden.
The naive reading was attractive: stronger cooperation had stabilized learning.
It was also underidentified.
A short path records displacement, not cause; it cannot distinguish efficient convergence from a missing update.
In a minimal matching task, hiding a binary partner type produced the opposite geometry---a path $44\times$ longer than with the type revealed---yet the hidden policy remained at chance.
The sign of the path-length change was environment-dependent while the learning failure was the same.

We trace this ambiguity to \textbf{relational aliasing}.
Let $r$ denote a relation that changes which action is useful: teammate identity, partner role, or task assignment.
When $r$ is hidden, trajectories from different relations update the same policy through the same observation.
Each conditional gradient may be informative, yet their mixture can be small because the gradients disagree.
This is not ordinary gradient noise: more samples estimate the same contracted expectation more accurately.
Nor is it the general statement that partial observability is hard.
It is a testable mechanism linking a specific information loss to the update vector.

We formalize the mechanism as \textbf{Information-Induced Gradient Contraction} (IIGC) and measure it with the \textbf{Relational Gradient Retention Ratio} $\kappa$.
The construction separates two questions that are otherwise conflated:
\emph{Is there relation-conditioned gradient energy?}
If so, \emph{what fraction survives aggregation?}
The first is answered by an energy gate $\mathcal{E}$; the second by $\kappa$.
Together with a behavioral path integral $\mathcal{P}$, they form a \textbf{Stable-or-Stuck diagnostic}: $\mathcal{P}$ describes motion, $\mathcal{E}$ tests whether a signal exists, and $\kappa$ tests whether mixing retains it.

The empirical design assigns a distinct role to each environment.
510K supplies the motivating ambiguity through checkpoint-level path geometry.
Toy Matching isolates relation-conditioned cancellation in a setting where the optimal actions are known.
Overcooked tests the diagnostic protocol at its boundary: sparse zero returns make the chosen REINFORCE estimator identically zero, so reporting a numerical $\kappa$ would be a category error.
Finally, 510K measurements under policy-gradient and TD update fields show that $\kappa$ is conditional on the estimator used to define $g_r$.

Our contributions are:
(1) an exact retained-versus-canceled energy decomposition for gradients mixed across hidden relations;
(2) a rollout-based estimator, $\kappa$, with explicit identifiability conditions and zero-energy handling;
(3) a diagnostic protocol that resolves path-length ambiguity without equating stillness with convergence; and
(4) controlled evidence across three environments, together with negative results that delimit what the diagnostic can claim.
The contribution is a measurement principle, not a new optimizer: before explaining a learning path, verify that a usable gradient survived the information bottleneck.


\section{Related Work}

\textbf{Partial observability and relational aliasing.}
Partially observable stochastic games and Dec-POMDPs formalize coordination with incomplete information~\cite{oliehoek2016concise,lowe2017multiagent,foerster2018counterfactual}.
The older notion of perceptual aliasing already identified the core representational problem: observationally identical situations may require different actions~\cite{chrisman1992perceptual}.
Recent diagnostics such as the $\lambda$-discrepancy test whether a learned state representation is Markovian~\cite{allen2024lambda}.
We study a narrower failure: the aliased variable is a relation, and the object of interest is the gradient energy lost when relation-conditioned updates share parameters.

\textbf{Gradient conflict and learning dynamics.}
Gradient variance and optimization geometry are central in supervised learning and policy optimization~\cite{bottou2018optimization,schulman2015trust,schulman2017proximal}.
Multi-task learning makes conflict explicit by comparing task gradients~\cite{yu2020gradient,zeng2021decentralized}; MARL adds credit assignment and non-stationarity~\cite{foerster2018counterfactual,hernandez2019emergent}.
Our ratio is related to gradient alignment but answers a different question: it measures the energy retained by the exact mixture induced by an unobserved relation and is reported only when the conditional energy is non-zero.

\textbf{Trajectory diagnostics and cooperation.}
Most empirical MARL comparisons emphasize final policies~\cite{papoudakis2019comparative}.
Behavioral archives, loss-landscape views, and checkpoint analyses instead study the route taken during learning~\cite{ecoffet2021first,li2018visualizing,lu2022discovered}.
Our path integral belongs to this family; our result is that geometry alone is non-identifying.
Self-play and league training establish strong learning systems~\cite{tesauro1995temporal,silver2016mastering,vinyals2019grandmaster,lanctot2017unified}, while social-dilemma and human--AI benchmarks vary cooperation structure~\cite{leibo2017multiagent,carroll2019utility}.
510K complements them with matched card mechanics and an intervention on team visibility.


\section{Gradient Retention under Hidden Relations}

\subsection{Setup and Behavioral Motion}

Let $r\in\{A,B\}$ be a relational configuration and let the policy receive $o_t=h(s_t)$ but not $r$.
The observation map is many-to-one whenever states from $A$ and $B$ share $o_t$ despite favoring different actions.
For a policy $\pi_\theta$ evaluated under relation $r$, define the return-weighted score-function gradient
\begin{equation}
\begin{split}
g_r &=
\mathbb{E}_{\tau\sim\pi_\theta,r}
\left[
\sum_t \nabla_\theta\log\pi_\theta(a_t\mid o_t)\,G_t
\right],\\[-2pt]
&\hspace{5.5cm} r\in\{A,B\},
\end{split}
\label{eq:grad}
\end{equation}
where $G_t$ is the discounted return-to-go.
We use a balanced mixture to expose the mechanism; unequal mixture weights replace the arithmetic mean by the corresponding weighted mean.

Behavioral motion is measured independently of the gradient.
For checkpoints $k=1,\ldots,K$, let $\mu_\pi^{(k)}\in\mathbb{R}^d$ be the mean feature vector over evaluation rollouts.
The path integral is
\begin{equation}
\mathcal{P}(\pi)=
\sum_{k=2}^{K}
\left\|\mu_\pi^{(k)}-\mu_\pi^{(k-1)}\right\|_2 .
\label{eq:path}
\end{equation}
It records total movement, not direction, return, or optimality.
Low $\mathcal{P}$ is therefore compatible with both convergence and paralysis; high $\mathcal{P}$ with both exploration and drift.

\medskip
\noindent\textbf{Proposition 0 (Path-Length Non-Identification).}
No decision rule based only on $\mathcal{P}$ can distinguish convergence from stagnation for all learning processes.
For $\mathcal{P}=0$, a policy initialized at an optimum and a policy with an identically zero update generate the same constant feature path.
For any $\mathcal{P}>0$, a directed path ending at an optimum and a loop returning to its initial behavior can be constructed with equal arc length.
The result is elementary but consequential: adding more precise path measurements cannot recover information that path geometry does not contain.
An independent signal about update direction or task value is necessary.

\subsection{Information-Induced Gradient Contraction}

With a balanced hidden relation, the expected update is
\begin{equation}
g_{\mathrm{mixed}}=\tfrac{1}{2}(g_A+g_B).
\label{eq:mixed}
\end{equation}
If $g_A$ and $g_B$ oppose one another, mixing erases signal that is present conditionally.
We call the loss of update energy caused by this many-to-one information map \textbf{Information-Induced Gradient Contraction}.

Define the total relation-conditioned energy
\begin{equation}
\mathcal{E}=\frac{\|g_A\|^2+\|g_B\|^2}{2}
\label{eq:energy}
\end{equation}
and, when $\mathcal{E}>0$, the \textbf{Relational Gradient Retention Ratio}
\begin{equation}
\kappa=
\frac{\left\|(g_A+g_B)/2\right\|^2}{\mathcal{E}}.
\label{eq:kappa}
\end{equation}
By convexity, $\kappa\in[0,1]$.
It is scale-free under a common rescaling of $g_A$ and $g_B$, but it retains both their angle and relative magnitudes.

\medskip
\noindent\textbf{Proposition 1 (Pythagorean Decomposition).}
Let $g_+=(g_A+g_B)/2$ and $g_-=(g_A-g_B)/2$.
Then
\begin{equation}
\mathcal{E}=\|g_+\|^2+\|g_-\|^2,
\qquad
\kappa=\frac{\|g_+\|^2}{\|g_+\|^2+\|g_-\|^2}.
\label{eq:pythagorean}
\end{equation}
Thus conditional gradient energy decomposes exactly into a retained component $\|g_+\|^2$ and a differential component $\|g_-\|^2$ removed by aggregation.
The identity follows by expanding $\|g_++g_-\|^2+\|g_+-g_-\|^2$; no distributional assumption is required.

\medskip
\noindent\textbf{Corollary 1 (Relation to Cosine Alignment).}
Writing $\theta_g$ for the angle between non-zero $g_A$ and $g_B$,
\begin{equation}
\kappa=
\frac{1}{2}
+
\frac{\|g_A\|\,\|g_B\|\cos\theta_g}
{\|g_A\|^2+\|g_B\|^2}.
\label{eq:cosine}
\end{equation}
When the norms match, $\kappa=(1+\cos\theta_g)/2$.
When one gradient is zero, $\kappa=1/2$ regardless of the direction of the other.
Accordingly, $1/2$ does not by itself mean orthogonality.

\medskip
\noindent\textbf{Proposition 2 (Effective Learning Signal).}
Under balanced mixing,
\begin{equation}
\|g_{\mathrm{mixed}}\|^2
=
\kappa\,\mathcal{E}.
\label{eq:effective}
\end{equation}
Small $\kappa$ implies contraction only when $\mathcal{E}$ is appreciable.
Conversely, $\mathcal{E}=0$ means that the chosen estimator contains no relation-conditioned signal to retain; $\kappa$ is then undefined, not zero.
This distinction separates cancellation from absence.

\medskip
\noindent\textbf{Extension to many relations.}
The binary, balanced case is expository rather than essential.
For relations $r=1,\ldots,m$ with mixture weights $p_r$, define
\begin{equation}
\bar g_p=\sum_{r=1}^{m}p_rg_r,
\qquad
\mathcal{E}_p=\sum_{r=1}^{m}p_r\|g_r\|^2.
\label{eq:weighted}
\end{equation}
A weighted variance identity gives
\begin{equation}
\mathcal{E}_p
=
\|\bar g_p\|^2
+
\sum_{r=1}^{m}p_r\|g_r-\bar g_p\|^2.
\label{eq:weighted-decomp}
\end{equation}
Thus $\kappa_p=\|\bar g_p\|^2/\mathcal{E}_p$ retains the same interpretation:
the first term is energy preserved by the actual training mixture, and the second is relation-specific energy removed by parameter sharing.
Unlike an unweighted pairwise cosine average, $\kappa_p$ respects relation frequency.
It also makes the intervention explicit: changing visibility may change each $g_r$, whereas rebalancing experience changes $p_r$.
These are different mechanisms and should not be conflated.

\medskip
\noindent\textbf{Analytic Toy case.}
Consider a Bernoulli policy with $\pi_\theta(a=1)=\sigma(\theta)$ and reward $+1$ for matching binary relation $r$ and $-1$ otherwise.
With $r=1$,
$g_1=2\sigma(\theta)(1-\sigma(\theta))$; with $r=0$,
$g_0=-2\sigma(\theta)(1-\sigma(\theta))$.
Each conditional gradient is non-zero away from saturation, but a balanced hidden mixture gives $\bar g=0$ and $\kappa=0$ for every $\theta$.
The example isolates the point of the ratio: the optimizer is not failing to estimate a direction; the information structure makes the population direction vanish.

\medskip
\noindent\textbf{Proposition 3 (Partition Refinement).}
Let $\mathcal{R}$ refine a coarser relation partition $\mathcal{C}$, and let
$g_c=\mathbb{E}[g_r\mid c]$.
Both partitions have the same mixed gradient $\bar g$, while Jensen's inequality gives
\begin{equation}
\sum_{r\in\mathcal{R}}p_r\|g_r\|^2
\ \geq\
\sum_{c\in\mathcal{C}}p_c\|g_c\|^2.
\label{eq:refinement}
\end{equation}
Whenever both denominators are positive,
$\kappa_{\mathcal{R}}\leq\kappa_{\mathcal{C}}$.
A coarse partition can therefore hide differential gradient energy by averaging it inside a relation label; refining the partition can reveal conflict but cannot manufacture additional retained energy.
This result makes relation choice part of the estimand.
``Teammate type,'' ``teammate identity,'' and ``private card'' are not interchangeable partitions even when they describe the same episode.

\subsection{Estimator Dependence}

The preceding identities hold for any two vectors, but their scientific meaning depends on how those vectors are defined.
For a policy-gradient estimator, $g_r^{\mathrm{PG}}$ changes action preferences.
A TD method instead supplies
\begin{equation}
\begin{split}
g_r^{\mathrm{TD}}
&=
\nabla_\theta
\mathbb{E}_{(s,a,r',s')\sim\mathcal{D}_r}
\left[
\bigl(Q_\theta(s,a)-y\bigr)^2
\right],\\
y&=r'+\gamma\max_{a'}Q_{\mathrm{target}}(s',a').
\end{split}
\label{eq:tdgrad}
\end{equation}

\medskip
\noindent\textbf{Proposition 4 (Gradient-Field Specificity).}
There is no estimator-independent ordering between
$\kappa(g_A^{\mathrm{PG}},g_B^{\mathrm{PG}})$ and
$\kappa(g_A^{\mathrm{TD}},g_B^{\mathrm{TD}})$.
The former compares action-preference updates; the latter compares prediction-error updates and also depends on replay distributions, features, and targets.
Changing the update field can change both angle and norm ratio even for the same environment and policy.
Hence cross-algorithm $\kappa$ values are meaningful only after naming the gradient estimator; a reversal is a property to explain, not evidence that one family is immune to hidden information.

\subsection{The Stable-or-Stuck Diagnostic}

The diagnostic is sequential.
First report $\mathcal{E}$.
If $\mathcal{E}$ is indistinguishable from zero, the estimator supplies no measurable signal and $\kappa$ is reported as N/A.
Otherwise report $\kappa$ and interpret it relative to a matched visible-relation baseline.
Finally use $\mathcal{P}$ to describe the resulting motion and return to check task success.

These are diagnostic descriptions, not universal classifiers: no environment-independent threshold on $\mathcal{P}$ or $\kappa$ is assumed.

\noindent\textbf{Diagnostic rule.}
After the energy gate, low $\mathcal{P}$/high $\kappa$ is compatible with convergence, low/low with contraction-induced stagnation, high/high with coherent motion, and high/low with conflict or drift; return determines whether any such geometry is useful.

\subsection{Finite-Sample Estimation Protocol}

For $N_r$ rollouts from each relation, we estimate
\begin{equation}
\widehat g_r
=
\frac{1}{N_r}\sum_{i=1}^{N_r}
\sum_t
\nabla_\theta\log\pi_\theta(a_{it}\mid o_{it})\,G_{it}.
\label{eq:estimator}
\end{equation}
The plug-in ratio $\widehat\kappa$ is not unbiased: both numerator and denominator are quadratic in noisy mean gradients.
The error is most consequential near the energy boundary, where a small denominator can turn sampling noise into an apparently decisive ratio.
Our reporting protocol therefore has five parts.

\begin{enumerate}
\item Freeze one checkpoint and hold policy parameters fixed across relations.
\item Intervene on $r$ while holding environment dynamics, reward, horizon, and rollout count fixed whenever the environment permits.
\item Report $\|\widehat g_r\|^2$ and $\widehat{\mathcal{E}}$, not only their ratio.
\item Obtain confidence intervals by resampling whole episodes within each relation; timestep resampling would break temporal dependence.
\item Report $\kappa$ only when the lower confidence bound on $\mathcal{E}$ clears a predeclared numerical floor.
\end{enumerate}

A useful falsification check splits rollouts from the \emph{same} relation into two batches.
Their retention estimates sampling stability, not relational contraction.
If split-batch agreement is low, the relation contrast is not yet estimable at the available rollout count.
If split-batch agreement is high but cross-relation $\kappa$ is low, cancellation rather than Monte Carlo instability is the sharper explanation.
This distinction is important for continuous or noisy relations, where two arbitrary rollout batches are not substitutes for two defined configurations.

\subsection{Three Distinct Estimands}

The phrase ``effect of hidden information'' can refer to three different quantities.
Separating them prevents a diagnostic result from silently becoming a causal or performance claim.

\textbf{Conditional conflict} is a property of one frozen policy:
\begin{equation}
\kappa(\pi;r)
=
\frac{\|\sum_r p_r g_r(\pi)\|^2}
{\sum_r p_r\|g_r(\pi)\|^2}.
\label{eq:conditional-estimand}
\end{equation}
It asks whether relation-specific updates agree at the current parameters.
Toy Matching identifies this quantity because partner type can be fixed during otherwise matched rollouts.

\textbf{Information effect} compares policies trained with and without observing $r$:
\begin{equation}
\Delta_\mathrm{info}
=
\mathbb{E}_{s}\!\left[\kappa(\pi^{\mathrm{visible}}_s;r)
-\kappa(\pi^{\mathrm{hidden}}_s;r)\right],
\label{eq:info-effect}
\end{equation}
where seed $s$, architecture, dynamics, reward, and checkpoints must be paired.
This estimand includes how information changes the learned representation and visitation distribution.
Comparing unrelated campaigns, input dimensions, or objectives does not identify it.

\textbf{Performance mediation} asks whether a change in return is caused through changed retention rather than another consequence of visibility.
That claim needs an intervention that changes $\kappa$ while controlling alternative paths, or a mediation design across matched seeds.
None of our benchmark results alone establishes mediation.
Accordingly, we use ``IIGC'' for measured loss of conditional update energy, reserve ``information effect'' for matched visibility contrasts, and report return as a separate outcome.


\section{Experimental Design}

We ask three questions.
\textbf{Q1: Can path geometry identify learning health?}
\textbf{Q2: When relation-conditioned gradient energy exists, does $\kappa$ expose cancellation?}
\textbf{Q3: What are the diagnostic's failure and transfer boundaries?}
The environments are assigned to these questions rather than treated as interchangeable replications.

\textbf{510K} is a four-player trick-taking game with 300 discrete actions and 112-dimensional observations.
Four modes share the base mechanics: \textsc{Single} is solo-competitive; \textsc{Static} uses fixed visible teams; \textsc{Dynamic} assigns teams through hidden red Aces; and \textsc{Obvious} uses \textsc{Dynamic}'s rules but appends a four-bit team indicator.
For Q1, 22 independent MaskablePPO runs~\cite{raffin2021stable,huang2022action} comprise 5 \textsc{Single}, 5 \textsc{Static}, 4 \textsc{Dynamic}, and 8 \textsc{Obvious} seeds.
We compute $\mathcal{P}$ from ten checkpoints at $10^5$-step intervals using seven normalized behavioral features (Appendix A).

\textbf{Toy Matching} has two actions and a balanced binary partner type.
The reward is $+1$ for matching the type and $-1$ otherwise.
\textsc{Revealed} includes a one-hot type; \textsc{Hidden} replaces it with zeros.
We train PPO for 10K steps over ten seeds per condition, then estimate relation-conditioned gradients from approximately 30 rollouts per type.

\textbf{Overcooked} pairs the learner with a chef or waiter in a sparse-reward kitchen.
\textsc{Static} reveals the partner role; \textsc{Dynamic} hides it and switches roles during an episode.
We train PPO for 1M steps over eight seeds per mode.
This experiment tests the energy gate because the return-weighted estimator in Eq.~\eqref{eq:grad} is exactly zero on zero-return rollouts.

For the estimator boundary, we additionally report 510K gradients from A2C (eight seeds), DQN (eight), and a discrete SAC variant (two).
These experiments are a separate campaign from the 22 path runs.
We report means and standard deviations across seeds; Toy path lengths are illustrative single-seed traces and are labeled as such.
Shared PPO hyperparameters are learning rate $3\times10^{-4}$, $\gamma=0.99$, GAE $\lambda=0.95$, batch size 64, and two 256-unit hidden layers.


\section{Results}

\subsection{Q1: Path Length Is Non-Identifying}

\begin{table}[t]
\centering
\small
\setlength{\tabcolsep}{3.5pt}
\begin{tabular}{lcccc}
\toprule
Mode & $n$ & $\mathcal{P}$ & Mean step & Step var. \\
\midrule
SINGLE  & 5 & $0.456\pm0.071$ & 0.051 & $6.1{\times}10^{-4}$ \\
STATIC  & 5 & $0.329\pm0.086$ & 0.037 & $3.1{\times}10^{-4}$ \\
OBVIOUS & 8 & $0.328\pm0.062$ & 0.036 & $2.3{\times}10^{-4}$ \\
DYNAMIC & 4 & $0.293\pm0.049$ & 0.033 & $2.8{\times}10^{-4}$ \\
\bottomrule
\end{tabular}
\caption{Behavioral path statistics for 22 independent 510K PPO runs.
The ordering is descriptive: modes differ in both objective and observability.}
\label{tab:paths}
\end{table}

\begin{figure}[t]
\centering
\includegraphics[width=\columnwidth]{figures/fig_path_lengths.pdf}
\caption{Behavioral path length across 510K modes.
\textsc{Dynamic}, with hidden team identity, has the shortest mean path; \textsc{Obvious}, with the same card hierarchy and visible teams, lies near \textsc{Static}.
Error bars show one standard deviation.}
\label{fig:paths}
\end{figure}

Table~\ref{tab:paths} and Figure~\ref{fig:paths} reproduce the motivating pattern:
$\textsc{Single}>\textsc{Static}\approx\textsc{Obvious}>\textsc{Dynamic}$ in mean path length.
The \textsc{Single}--\textsc{Dynamic} standardized difference is large ($d=2.67$), and step-size variance also falls by $2.2\times$.
This establishes less behavioral motion, not better learning.
The modes change reward and termination rules as well as information, so the four-way ordering is not a causal dose-response curve.
The matched-mechanics contrast, \textsc{Obvious} versus \textsc{Dynamic}, is directionally consistent with an information effect ($0.328$ versus $0.293$) but is underpowered and unpaired.
The correct conclusion from 510K alone is therefore ambiguity: hidden relations can coexist with a short, low-variance path.

\subsection{Q2: Toy Matching Identifies Contraction}

\begin{table}[t]
\centering
\small
\setlength{\tabcolsep}{5pt}
\begin{tabular}{lcccc}
\toprule
Condition & $\mathcal{P}$ & $\kappa$ & Return & $n$ \\
\midrule
REVEALED & 0.011 & $0.726\pm0.100$ & $19.2\pm0.3$ & 10 \\
HIDDEN   & 0.480 & $0.039\pm0.073$ & $-0.3\pm0.3$ & 10 \\
\bottomrule
\end{tabular}
\caption{Toy PPO results.
Path lengths are illustrative single-seed traces; $\kappa$ and return aggregate ten independent seeds.}
\label{tab:toy}
\end{table}

\begin{figure}[t]
\centering
\includegraphics[width=\columnwidth]{figures/kappa_figure.pdf}
\caption{Toy $\kappa$ by seed.
Revealing partner type preserves most conditional gradient energy; hiding it drives the retained fraction toward zero.
Dashed lines denote condition means.}
\label{fig:kappa}
\end{figure}

Toy Matching reverses the 510K geometry.
The hidden policy travels $44\times$ farther in the illustrative trace, yet its return stays at chance.
Thus neither a short nor a long path has a fixed implication for learning health.
The conditional gradients supply the missing identification:
revealing type yields $\kappa=0.726\pm0.100$, while aliasing the two types yields $0.039\pm0.073$.
Seven of ten hidden seeds are below $0.005$; all revealed seeds exceed $0.54$.
The near-zero retained fraction and chance return agree even though $\mathcal{P}$ points in the opposite direction.

This experiment also distinguishes contraction from sampling variance.
The two partner types have opposite optimal actions by construction, and each conditional rollout distribution supplies a non-zero learning signal.
Increasing rollout count can reduce estimation error around $g_A$ and $g_B$, but it cannot change the fact that their balanced expectation cancels when type is unobserved.

\subsection{Q3: The Zero-Energy Boundary}

\begin{table}[t]
\centering
\small
\setlength{\tabcolsep}{3.5pt}
\resizebox{\columnwidth}{!}{%
\begin{tabular}{lcccc}
\toprule
Mode & Reported $\kappa$ & Chef ret. & Waiter ret. & Energy status \\
\midrule
STATIC  & $0.497\pm0.006$ & 186.8 & 43.5 & One-sided in 4/8 \\
DYNAMIC & N/A & 0.0 & 0.0 & Zero in 8/8 \\
\bottomrule
\end{tabular}}
\caption{Overcooked PPO stress test over eight seeds.
Zero return makes the return-weighted gradients, and hence $\mathcal{E}$, zero in \textsc{Dynamic}; $\kappa$ is undefined.
Static averages include four seeds with substantial return under both partner roles.}
\label{tab:overcooked}
\end{table}

Overcooked does not provide evidence of cancellation; it provides evidence for the energy gate.
Every \textsc{Dynamic} seed has zero return under both frozen partner evaluations.
Equation~\eqref{eq:grad} then gives $g_A=g_B=0$ and $\mathcal{E}=0$.
A numerical implementation with a clamped denominator returns a sentinel near zero, but interpreting that value as $\kappa\approx0$ would confuse absent gradient energy with opposing gradients.
We report N/A instead.

The \textsc{Static} mean, $0.497$, requires equal care.
Four seeds obtain substantial return under both partner roles; the other four succeed almost exclusively with the chef.
When one conditional gradient dominates, Corollary~1 forces $\kappa$ toward $1/2$ even without orthogonality.
Overcooked therefore supports a protocol-level conclusion: report conditional energy and return alongside the ratio.
It does not support the stronger claim that hidden roles caused measured cancellation.

\subsection{Objective-Dependent Retention}
\label{sec:cross}

\begin{table}[t]
\centering
\small
\setlength{\tabcolsep}{3.5pt}
\begin{tabular}{llccc}
\toprule
Estimator & Update field & $\kappa_{\rm SINGLE}$ & $\kappa_{\rm DYNAMIC}$ & $n$ \\
\midrule
A2C & Policy grad. & $0.644\pm0.201$ & $0.519\pm0.060$ & 8 \\
DQN & TD loss      & $0.797\pm0.123$ & $0.917\pm0.063$ & 8 \\
SAC & Hybrid       & $0.542\pm0.060$ & $0.569\pm0.017$ & 2 \\
\bottomrule
\end{tabular}
\caption{510K retention under different update fields.
Rows are not a common algorithm score: each defines $g_r$ through a different objective.
SAC is exploratory ($n=2$).}
\label{tab:cross}
\end{table}

Table~\ref{tab:cross} shows why estimator naming is substantive.
For A2C policy gradients, hiding team identity lowers the retained fraction by $0.125$.
For DQN's TD-loss field, the direction reverses by $0.120$; the two-seed SAC result is also slightly higher under \textsc{Dynamic}.
Proposition~3 precludes reading this as policy-gradient vulnerability versus value-based immunity.
The TD vectors reflect prediction residuals under a replay distribution, not the action-preference mixture in Eq.~\eqref{eq:grad}.
The reversal instead establishes a boundary: $\kappa$ characterizes a specified update field.

We exclude the available 510K PPO contrast from this table because the visible side contains only one recovered run while \textsc{Dynamic} contains nine.
Vanilla REINFORCE is similarly inconclusive: its across-seed standard deviation exceeds $0.23$ in both modes and several runs have zero or one-sided gradient energy.
These negative results do not vanish under better prose; they limit the algorithm-level claim to the matched A2C and DQN measurements.

\subsection{Robustness and Claim Audit}

Three checks separate stable evidence from attractive interpretation.
First, the 510K path ordering does not depend on one hand-designed coordinate:
deleting any of the seven behavioral features preserves
$\textsc{Single}>\textsc{Static}>\textsc{Dynamic}$ in all seven leave-one-feature-out analyses.
This supports robustness to feature omission, but not invariance to a different representation.

Second, the Toy contrast is not driven by its mean.
The smallest revealed PPO value is $0.547$; the largest hidden value is $0.235$.
An independent A2C campaign repeats the direction
($0.842\pm0.025$ revealed versus $0.195\pm0.352$ hidden, ten seeds each), but two hidden seeds retain high energy and make the variance large.
We therefore use PPO for the quantitative claim and A2C only as a directional replication.

Third, every ratio is checked against the update that produced it.
The Overcooked zeros fail the energy gate.
The DQN values pass a different scientific test than the A2C values because they use TD-loss gradients.
The unmatched PPO baseline, two-seed SAC sweep, and same-condition reveal
audit in Figure~\ref{fig:split-batch} remain exploratory.

\begin{figure}[t]
\centering
\includegraphics[width=\columnwidth]{figures/fig_reveal.pdf}
\caption{Split-batch retention across five 510K reveal conditions
($n=2$; bars show one standard deviation).
Same-condition batches audit sampling consistency, not relation-conditioned
retention or a causal information effect.}
\label{fig:split-batch}
\end{figure}

Table~\ref{tab:evidence} records this claim hierarchy.

\begin{table}[t]
\centering
\small
\setlength{\tabcolsep}{2.5pt}
\begin{tabular}{
>{\raggedright\arraybackslash}p{0.20\columnwidth}
>{\raggedright\arraybackslash}p{0.34\columnwidth}
>{\raggedright\arraybackslash}p{0.34\columnwidth}}
\toprule
Evidence & Supports & Does not support \\
\midrule
510K paths &
Mode-dependent motion &
Convergence or causal mechanism \\
Toy PPO &
Cancellation under controlled aliasing &
Benchmark-wide prevalence \\
Overcooked &
Need for the energy gate &
Measured cancellation \\
510K A2C/DQN &
Estimator-dependent retention &
Universal algorithm ranking \\
\bottomrule
\end{tabular}
\caption{The inferential role of each experiment.
No result is asked to establish both mechanism and generality.}
\label{tab:evidence}
\end{table}

\subsection{What Would Falsify IIGC?}

The framework is falsifiable at three levels.
At the algebraic level, if the measured mixed update is not the weighted sum of the reported conditional updates, the relation partition or gradient implementation is wrong.
At the estimator level, if cross-relation $\kappa$ remains low after revealing $r$ to an expressive policy and split-batch agreement is equally low, sampling error is a more parsimonious account than information-induced contraction.
At the task level, if low $\kappa$ coexists with high return because a compromise policy is optimal under the mixture, contraction is real but not harmful.
IIGC predicts lost update energy; it does not assert that every lost component is needed for the mixture objective.

The last case matters in cooperative games.
Two relation-specific optima may disagree while a shared policy is still adequate.
Conversely, high $\kappa$ can accompany low return when both conditional gradients align toward a poor local solution.
Return, energy, retention, and path length answer different questions.
A valid empirical story preserves those distinctions rather than forcing them into a single score.


\section{Discussion}

\textbf{What the diagnostic identifies.}
A path integral answers ``how much did behavior move?''
The energy gate answers ``did the chosen estimator contain a conditional signal?''
Only after that gate does $\kappa$ answer ``how much survived relation mixing?''
This ordering prevents two common errors: calling a zero-gradient run cancellation, and calling a short path convergence.
Toy demonstrates the intended use because the conditional energy, optimal actions, and return are all identifiable.
510K demonstrates why the question arises.
Overcooked demonstrates when the ratio must be withheld.

\textbf{Use and intervention.}
The diagnostic requires rollout access under identifiable relations, not gradients saved during training.
It can therefore be applied to a checkpoint before a full training campaign.
A low $\kappa$ relative to a matched visible-relation baseline motivates an information intervention---for example, exposing team identity or learning a memory state---but does not prove that any particular repair will improve return.
Likewise, cross-estimator comparisons may motivate algorithm choice only after the update fields and energy scales are made comparable.
The measurement is a mechanism check, not an early-stopping oracle.

\textbf{Limits.}
First, $\kappa$ requires relation-conditioned rollouts; a continuously valued or unobserved relation must be identified or discretized before measurement.
Second, the balanced two-relation analysis does not cover arbitrary mixture weights, although the energy identity extends by replacing the mean with a weighted projection.
Third, the 510K path integral depends on seven domain features; leave-one-feature-out tests preserve the reported ordering, but a representation-independent path metric remains open.
Fourth, the empirical campaigns are uneven: the 510K path and gradient studies use different checkpoints, the PPO visible baseline is incomplete, SAC has two seeds, and Overcooked lacks a comparable behavioral path.
Fifth, $\kappa$ depends on the estimator, reward scale through relative gradient magnitudes, and rollout distribution.
Broader tests on standard hidden-information benchmarks such as Hanabi or SMAC are needed before claiming generality.

\textbf{Parameterization and metric choice.}
The Euclidean norm in Eq.~\eqref{eq:kappa} is not invariant to a nonlinear reparameterization of the policy network.
The comparison is therefore valid across relations for a fixed architecture and checkpoint, but raw values should not be compared across parameterizations as though they lived in a common geometry.
For any fixed positive semidefinite metric $M$, the construction generalizes to
\begin{equation}
\mathcal{E}_M=\sum_r p_r g_r^\top M g_r,
\qquad
\kappa_M=\frac{\bar g^\top M\bar g}{\mathcal{E}_M}.
\label{eq:metric}
\end{equation}
The weighted decomposition in Eq.~\eqref{eq:weighted-decomp} still holds in the $M$-seminorm.
Choosing $M$ from the inverse Fisher information would measure retention in a local policy-distribution geometry rather than raw parameter coordinates, but estimating that metric reliably is a separate problem.
We use Euclidean energy and keep architecture fixed; we do not claim reparameterization invariance.

\textbf{Estimator controls.}
A common positive scaling of both conditional gradients cancels in $\kappa$, but relation-specific reward normalization does not.
Rollout horizon, discount, advantage construction, action masks, and baseline must therefore match across $r$.
An action-independent control variate leaves the population policy gradient unchanged, yet can change finite-sample $\widehat\kappa$ by changing covariance.
For TD gradients, replay composition and target-network age are equally part of the measurement.
These are not implementation footnotes: they specify which vector field the ratio describes.

\begin{table}[t]
\centering
\small
\setlength{\tabcolsep}{2.5pt}
\begin{tabular}{
>{\raggedright\arraybackslash}p{0.19\columnwidth}
>{\raggedright\arraybackslash}p{0.35\columnwidth}
>{\raggedright\arraybackslash}p{0.34\columnwidth}}
\toprule
Observed pattern & Next measurement & Not yet justified \\
\midrule
Zero energy &
Audit reward, coverage, and estimator &
Calling the failure cancellation \\
Unstable split &
Add rollouts or variance reduction &
Comparing relations \\
Low cross-$r$ $\kappa$ &
Reveal $r$ or add memory in a matched run &
Claiming return will improve \\
High $\kappa$; low return &
Audit representation, credit, and objective &
Ruling out learning failure \\
Reversal &
Match data and update fields &
Ranking algorithm families \\
\bottomrule
\end{tabular}
\caption{Operational use of the diagnostic.
Each pattern determines a next measurement, not an automatic intervention.}
\label{tab:actions}
\end{table}

\textbf{Scope of the claim.}
The framework does not assert that hidden information always lowers $\kappa$, that low $\kappa$ always lowers return, or that a particular optimizer repairs contraction.
It establishes an exact accounting identity, a conditional measurement protocol, and one controlled case where relational aliasing erases an otherwise non-zero update.
The benchmark studies probe transfer and expose failure conditions.
This narrower statement is also more useful: it says which quantity can be measured, which controls make the comparison interpretable, and which result would falsify the mechanism.

\textbf{Design implications.}
The first intervention depends on which gate fails.
When $\mathcal{E}\approx0$, changing observation visibility may be premature: the return estimator, reward density, or state coverage supplies no conditional direction to compare.
When energy is present and split-batch agreement is high but cross-relation $\kappa$ is low, the conflict is localized to the relation partition.
Then exposing $r$, adding memory, or conditioning the policy on an inferred relation is a targeted test.
The aim is not to maximize $\kappa$ in isolation, but to restore a direction that improves the mixture objective.

Changing the optimizer answers a different question.
Clipping, replay, advantage estimation, or a value objective can rotate and rescale $g_r$; a higher resulting ratio may reflect a different update field rather than resolution of the original conflict.
A fair algorithm comparison must therefore report both fields on matched data or evaluate a common probe gradient after training.
Without that control, $\kappa$ explains the optimizer's update geometry but does not rank its competence.

Finally, visibility itself is not binary.
Proposition~3 implies that a coarse label may average away conflict that appears under identities or episode-level assignments.
A practical study should predeclare the finest relation partition it can intervene on, then report coarser aggregations as sensitivity analyses.
This turns an informal statement such as ``hidden teamwork is hard'' into a sequence of auditable questions: which relation, under which mixture, for which update field, at what energy, and with what return?

\textbf{Next steps.}
The immediate empirical need is not another environment count but a matched intervention: identical dynamics, reward, seeds, and checkpoints, changing only whether $r$ is observed.
The theoretical extension is an $m$-relation decomposition with explicit mixture weights and confidence intervals for both $\mathcal{E}$ and $\kappa$.
Online estimates could then test whether contraction predicts the benefit of memory, teammate inference, or information reveal before those mechanisms are added.


\section{Conclusion}

Behavioral stillness is not evidence of convergence.
Under hidden relations, useful conditional gradients can be destroyed by aggregation; in sparse-reward settings, the conditional gradients may instead be absent.
The energy $\mathcal{E}$ and retention ratio $\kappa$ separate these cases, while the path integral $\mathcal{P}$ describes their visible consequence.
The resulting rule is simple: establish that a signal exists, measure what mixing retains, and only then interpret the learning path.


\clearpage
\bibliography{references}

\clearpage
\section*{Appendix A: 510K Game Rules and Feature Definitions}

\textbf{Game rules.}
510K is a four-player trick-taking game with a standard 52-card deck.
Cards 5, 10, and K carry point values (5, 10, and 10).
The lead player sets the pattern (single, pair, straight, bomb, or 510K pattern); others must follow with a higher-rank play or pass.
The trick winner collects all 5/10/K points from the trick.

Four modes share base card mechanics but differ in team structure and termination:
\begin{itemize}
\item \textsc{Single}: solo-competitive. The game ends when the first player empties its hand. The winner receives $+30$ and each loser $-10$.
\item \textsc{Static}: fixed teams (0\&2 versus 1\&3). The game ends when both members of one team finish. Winners receive $+15$ and losers $-15$.
\item \textsc{Dynamic}: teams follow the hidden red-Ace distribution (holders of A$\heartsuit$/A$\diamondsuit$ form one team). Scoring matches \textsc{Static}; red A also changes the card hierarchy.
\item \textsc{Obvious}: identical to \textsc{Dynamic} in card mechanics, with team membership added to the observation as a four-bit indicator.
\end{itemize}

\textbf{Behavioral features.}
Seven mode-independent features ($\phi_1$--$\phi_7$), all with VIF $<2$, define the 510K path:
\begin{itemize}
\item $\phi_1$ MyScore: cumulative 510K points, normalized by 150.
\item $\phi_2$ MyHandSize: fraction of the original hand played.
\item $\phi_3$ MyStrength: average card rank mapped to $[0,1]$.
\item $\phi_4$ TrickScore: points at stake in the current trick, normalized by 100.
\item $\phi_5$ PassCount: consecutive passes in the current trick, normalized.
\item $\phi_6$ ScoreSpread: standard deviation of player scores, normalized by 50.
\item $\phi_7$ SuppressionGap: rank gap to the leader's strongest card, normalized by 12.
\end{itemize}
Deleting any single feature preserves the mean ordering
$\textsc{Single}>\textsc{Static}>\textsc{Dynamic}$ in all seven leave-one-feature-out checks.

\textbf{Observation intervention.}
\textsc{Obvious} has a 116-dimensional observation (112 base features plus four team-indicator bits).
Bit $i$ is one when player $i$ is on the observing agent's team.
\textsc{Dynamic} omits these four bits.

\end{document}
